\chapter{مبرهنة سيغل--فالفش والتوزيع المنتظم للأعداد الأولية}

\section{نطاق الفصل وحالته}

\noindent\textbf{حالة الفصل:} \textenglish{REVIEWED}. يرقّي هذا الفصل نتيجة
الفصل العاشر من مبرهنة نوعية عند ترديد ثابت إلى تقدير كمي موحد عندما ينمو
الترديد مع \(x\) في المجال
\[
q\le (\log x)^A.
\]
والنتيجة المركزية هي أنه لكل \(A>0\) يوجد ثابت \(c_A>0\) بحيث، بانتظام في
الفئات المختزلة \(a\pmod q\)،
\[
\psi(x;q,a)
=
\frac{x}{\varphi(q)}
+O_A\!\left(xe^{-c_A\sqrt{\log x}}\right).
\]
الثابت \(c_A\)، والثوابت الضمنية الناتجة من المسار غير المشروط العام، غير
فعالة بسبب استعمال مبرهنة سيغل لإبعاد الصفر الاستثنائي.

بدأ التأليف بعد إغلاق بوابة
\textenglish{EVIDENCE-FIRST / PRE-AUTHORING}، ثم اجتاز الفصل تدقيق ما بعد
التأليف بحكم \textenglish{PASS} ونجح في فحص الجودة وبناء PDF الكامل.
وتوجد ملفات البوابة وتدقيق التأليف وخريطة البرهان والديون التحليلية في
قائمة الوثائق الآتية:
\begin{itemize}
  \item سجل الأدلة للفصل 12 في مجلد مراجعات الأدبيات.
  \item خريطة البرهان للفصل 12 في مجلد مراجعات الأدبيات.
  \item تدقيق بيرون وتحويل المسار والصيغة الصريحة.
  \item تدقيق الصفر الاستثنائي وعدم فعالية الثوابت.
  \item التدقيق المنطقي للفصل.
  \item التحقق المرجعي للفصل.
  \item تقرير المراجعة المستقلة.
  \item وثيقة إغلاق المراجعة.
\end{itemize}

لا يثبت الفصل صيغة بيرون المقطوعة أو تحويل المسار من الصفر. بل يقتبس صيغة
صريحة موحدة محددة، ثم يثبت داخليًا اختيار ارتفاع القطع، وضبط الجزء غير
الاستثنائي، وامتصاص الصفر الاستثنائي، وتجميع الشخصيات. ولا يستعمل الفصل
مبرهنة Bombieri--Vinogradov أو Linnik أو فرضية ريمان المعممة. وقد اجتاز
الفصل مراجعة ثانية مستقلة بحكم \textenglish{APPROVED}، ورُقّي إلى
\textenglish{REVIEWED} ودُمج عبر PR \#20. ولا تعني هذه الحالة أنه أصبح
\textenglish{RELEASE-READY}.

\subsection{ثلاث طبقات يجب عدم خلطها}

\begin{enumerate}
  \item يثبت الفصل العاشر أنه لكل \(q\) ثابت:
  \[
  \psi(x;q,a)\sim \frac{x}{\varphi(q)}.
  \]
  لا تسمح هذه العبارة باستبدال \(q\) بترديد يعتمد على \(x\).
  \item يثبت هذا الفصل الانتظام الفردي لكل
  \(q\le(\log x)^A\)، مع حد خطأ موحد وغير فعال.
  \item مبرهنة Bombieri--Vinogradov نتيجة متوسطية على الترديدات في مجال
  أكبر بكثير، وتحتاج أدوات إضافية؛ لذلك تبقى مؤجلة.
\end{enumerate}

\subsection{نواتج التعلم}

بعد إتمام الفصل ينبغي أن يستطيع القارئ:
\begin{enumerate}
  \item تفسير الفرق بين الترديد الثابت والانتظام في الترديد.
  \item رد مجموع فون مانغولت لشخصية مستحثة إلى جدها البدائي.
  \item استعمال الصيغة الصريحة المقطوعة من دون إخفاء ديون بيرون وتحويل المسار.
  \item اختيار ارتفاع قطع من رتبة \(e^{\kappa\sqrt{\log x}}\).
  \item عزل حد الصفر الاستثنائي قبل أخذ القيم المطلقة.
  \item تحديد الموضع الوحيد الذي تستعمل فيه مبرهنة سيغل.
  \item اشتقاق صيغ موحدة للدوال \(\psi\)، و\(\vartheta\)، و\(\pi\).
  \item تفسير سبب عدم فعالية الثوابت النهائية.
\end{enumerate}

\section{دوال العد والمرشح بالشخصيات}

لكل شخصية ديريشليه \(\chi\pmod q\)، نضع
\[
\psi(x,\chi)=\sum_{n\le x}\chi(n)\Lambda(n).
\]
ولكل فئة مختزلة \((a,q)=1\)، نعيد استعمال تعامد الشخصيات لنحصل على الهوية
الدقيقة
\[
\psi(x;q,a)
=
\frac1{\varphi(q)}
\sum_{\chi\bmod q}\overline{\chi(a)}\,\psi(x,\chi).
\tag{12.1}
\]
هذه هي النسخة المنتهية من التفكيك المسجل في
\textenglish{\texttt{ANT-PROP-10-01}}. ولأن عدد الشخصيات هو \(\varphi(q)\)،
فإن تقديرًا موحدًا لكل شخصية لا يخسر عاملًا إضافيًا من رتبة \(q\) عند
التجميع في \((12.1)\).

\section{الرد إلى الجد البدائي}

\begin{lemma}[رد مجموع فون مانغولت إلى الشخصية البدائية]
\label{lem:chapter12-primitive-psi-reduction}
\resultid{ANT-LEM-12-01}
\provedhere
لتكن \(\chi\pmod q\) مستحثة من شخصية بدائية \(\chi^*\pmod r\)، حيث
\(r\mid q\). عندئذ، لكل \(x\ge2\)،
\[
\psi(x,\chi)
=
\psi(x,\chi^*)+O(\log(2q)\log(2x)),
\]
بثابت مطلق.
\end{lemma}

\begin{proof}
إذا كان \((n,q)=1\)، فإن \(\chi(n)=\chi^*(n)\). ولا يمكن أن يختلف الطرفان
على قوة أولية \(p^k\) إلا إذا كان \(p\mid q\) وحذف العامل المحلي الموافق
عند الاستحثاث. لذلك
\[
\psi(x,\chi)-\psi(x,\chi^*)
=
-\sum_{\substack{p^k\le x\\p\mid q,\ p\nmid r}}
\chi^*(p)^k\log p.
\]
ومن ثم
\begin{align*}
\left|\psi(x,\chi)-\psi(x,\chi^*)\right|
&\le
\sum_{p\mid q}\left\lfloor\frac{\log x}{\log p}\right\rfloor\log p\\
&\le \omega(q)\log x\\
&\ll \log(2q)\log(2x).
\end{align*}
\end{proof}

في المجال \(q\le(\log x)^A\)، يكون هذا الخطأ متعدد اللوغاريتمات، ولذلك
يمتص في كل من الحد الأسي المستهدف وأي ادخار لوغاريتمي ثابت.

\section{مدخلان مقتبسان من النظرية الكمية}

المسار الكمي لا ينتج من مبرهنة Wiener--Ikehara النوعية وحدها. نحتاج حدًا
فعالًا للشخصية الرئيسية، وصيغة صريحة مقطوعة موحدة للشخصيات البدائية.

\begin{theorem}[حد دو لا فاليه بوسان لمبرهنة الأعداد الأولية]
\label{thm:chapter12-effective-pnt}
\resultid{ANT-THM-12-01}
\citedresult{\cite{DeLaValleePoussin1896,Davenport2000,MontgomeryVaughan2007}}
توجد ثوابت مطلقة فعالة \(c,C>0\) بحيث، لكل \(x\ge3\)،
\[
\psi(x)=x+O\!\left(Cxe^{-c\sqrt{\log x}}\right).
\tag{12.2}
\]
\end{theorem}

هذه النتيجة أقوى كميًا من \textenglish{\texttt{ANT-THM-09-03}}، الذي يثبت
فقط \(\psi(x)\sim x\). لا نستخرج \((12.2)\) من النتيجة النوعية.

لتجنب مشكلة نقاط القفز، نعرف النسخة المنصفة
\[
\psi_0(x,\chi)
=
\frac12\left(\sum_{n<x}\chi(n)\Lambda(n)+
              \sum_{n\le x}\chi(n)\Lambda(n)\right).
\]
ولدينا دائمًا
\[
\psi(x,\chi)-\psi_0(x,\chi)=O(\log(2x)).
\]

\begin{theorem}[مدخل كمي مركب للصيغة الصريحة بعد عزل الاستثناء]
\label{thm:chapter12-truncated-explicit-formula}
\resultid{ANT-THM-12-02}
\citedresult{\cite{Davenport2000,MontgomeryVaughan2007}}
توجد ثوابت مطلقة فعالة \(c,C>0\) بالخاصية الآتية. لتكن \(\chi\) شخصية
بدائية غير رئيسية بموصل \(r\)، ولتكن \(x\ge3\) و\(r\le T\le x\). يمكن
اختيار ارتفاع
\[
U\in[T,2T]
\]
لا يساوي ارتفاع صفر، بحيث
\[
\psi_0(x,\chi)
=
-\mathbf 1_{\beta_\chi}
\frac{x^{\beta_\chi}}{\beta_\chi}
+O\!\left(
xe^{-\frac{c\log x}{\log(r(U+2))}}\log^2(r(U+2))
+
\frac{x\log^2(rxU)}{U}
+
\log^2(rxU)
\right).
\tag{12.3}
\]
هنا \(\beta_\chi\) هو الصفر الحقيقي الاستثنائي الممكن في المنطقة القياسية
من الفصل الحادي عشر، ويكون الحد الأول غائبًا إذا لم يوجد هذا الصفر.
الثوابت في \((12.3)\) فعالة قبل امتصاص الحد الاستثنائي.
\end{theorem}

الصيغة \((12.3)\) ليست نقلًا حرفيًا لمبرهنة مفردة، بل مدخلًا مركبًا من:
الصيغة الصريحة لـ\(\psi(x,\chi)\)، واختيار ارتفاع قطع جيد، وتقدير عد
الأصفار، والمنطقة القياسية الخالية. مواضع المسار الكلاسيكي هي Davenport،
الطبعة الثانية، الفصول 20--23، الصفحات 115--134، وMontgomery--Vaughan،
الفصلان 11--12، الصفحات 358--418. وقد ضمت مساهمات الأصفار البديهية وعامل
غاما وخطأ القطع في الحد المكتوب. لا نسجل هذه النتيجة
\textenglish{PROVED-HERE} لأن إثباتها الكامل يحتاج صيغة بيرون المقطوعة
وتحويل المسار وتقديرات رأسية لم تغلق بعد داخل الموسوعة.

\section{اختيار ارتفاع القطع وضبط الجزء غير الاستثنائي}

\begin{lemma}[اختيار الارتفاع اللوغاريتمي]
\label{lem:chapter12-nonexceptional-zero-bound}
\resultid{ANT-LEM-12-02}
\provedhere
ليكن \(A>0\). توجد ثوابت \(\kappa_A,c_A,C_A>0\) فعالة بحيث إذا
\[
q\le(\log x)^A,
\qquad
\chi^*\pmod r,
\qquad
r\mid q,
\]
وكانت \(\chi^*\) بدائية وغير رئيسية، فإن الجزء غير الاستثنائي في
\((12.3)\) يحقق
\[
\psi(x,\chi^*)
+
\mathbf 1_{\beta_{\chi^*}}
\frac{x^{\beta_{\chi^*}}}{\beta_{\chi^*}}
\ll_A
xe^{-c_A\sqrt{\log x}}.
\tag{12.4}
\]
الثوابت هنا فعالة؛ عدم الفعالية لم تدخل بعد.
\end{lemma}

\begin{proof}
ضع \(L=\log x\)، واختر
\[
T=e^{\kappa\sqrt L},
\]
حيث \(\kappa>0\) ثابت صغير سيحدد بحسب الثوابت المطلقة في
\((12.3)\). عند كبر \(x\)، لدينا
\[
r\le q\le L^A\le T\le x.
\]
طبّق المبرهنة \ref{thm:chapter12-truncated-explicit-formula} واختر الارتفاع
\(U\in[T,2T]\) الذي تمنحه. عندئذ
\[
\log(r(U+2))
\le A\log L+\kappa\sqrt L+O(1)
\ll_{A,\kappa}\sqrt L.
\]
ومن ثم
\[
\exp\!\left(-\frac{cL}{\log(r(U+2))}\right)
\le e^{-c_{A,\kappa}\sqrt L}.
\]
وتعطي كلفة القطع، باستعمال \(U\ge T\) و\(U\le2T\)،
\[
\frac{x\log^2(rxU)}{U}
\ll_A
xe^{-\kappa\sqrt L}L^2,
\]
بينما لا يتجاوز الحد الأخير قوة ثابتة من \(L\). وتمتص كل خسارة متعددة
اللوغاريتمات بعد تصغير ثابت الأس، لأن
\[
L^D e^{-c\sqrt L}\le e^{-c'\sqrt L}
\]
لكل \(D\) ثابت ولكبر \(L\). وأخيرًا يمتص الفرق بين \(\psi\) و\(\psi_0\)
في الحد نفسه. أما المجال المحدود من قيم \(x\)، فيمتص بتكبير الثابت
\(C_A\).
\end{proof}

\section{الصفر الاستثنائي وموضع استعمال مبرهنة سيغل}

اللمّة السابقة لا تسمح بحذف الحد
\[
-\frac{x^{\beta_\chi}}{\beta_\chi}.
\]
يجب إبقاؤه ظاهرًا حتى هذه النقطة.

\begin{lemma}[امتصاص حد الصفر الاستثنائي]
\label{lem:chapter12-exceptional-zero-absorption}
\resultid{ANT-LEM-12-03}
\provedhere
ليكن \(A>0\). يوجد ثابت غير فعال \(c_A>0\) بحيث إذا
\(r\le q\le(\log x)^A\)، وكانت شخصية حقيقية بدائية بموصل \(r\) تملك
صفرًا استثنائيًا \(\beta\)، فإن
\[
\frac{x^\beta}{\beta}
\ll_A
xe^{-c_A\sqrt{\log x}}.
\tag{12.5}
\]
\end{lemma}

\begin{proof}
الصفر الاستثنائي قريب من الواحد، ولذلك \(\beta\ge1/2\) بعد امتصاص مجال
محدود من الموصلات، ومن ثم \(1/\beta\le2\). من
\textenglish{\texttt{ANT-COR-11-01}}، لكل \(\varepsilon>0\)،
\[
1-\beta\ge c_\varepsilon r^{-\varepsilon},
\]
حيث \(c_\varepsilon\) غير فعال. اختر \(\varepsilon>0\) بحيث
\(A\varepsilon=1/2\). عندئذ
\[
r^{-\varepsilon}
\ge q^{-\varepsilon}
\ge(\log x)^{-A\varepsilon}
=(\log x)^{-1/2}.
\]
إذًا
\[
(1-\beta)\log x
\ge c_\varepsilon\sqrt{\log x},
\]
ومن ثم
\[
x^\beta
=x\exp(-(1-\beta)\log x)
\le xe^{-c_\varepsilon\sqrt{\log x}}.
\]
هذا هو الموضع الوحيد الذي تدخل فيه مبرهنة سيغل في المسار الأدنى، ويرث
الثابت عدم فعاليتها.
\end{proof}

\begin{remark}[لماذا لا يجوز إخفاء الاستثناء مبكرًا؟]
قبل اللمّة السابقة قد يكون \(x^\beta\) قريبًا من \(x\). المنطقة القياسية
وحدها لا تبعد \(\beta\) بما يكفي عندما ينمو الموصل. لذلك فإن وضع الحد في
رمز \(O\) قبل استعمال سيغل يفقد المعلومة الحاسمة ويخفي مصدر عدم الفعالية.
\end{remark}

\section{مبرهنة سيغل--فالفش}

\begin{theorem}[مبرهنة سيغل--فالفش في صيغة \(\psi\)]
\label{thm:chapter12-siegel-walfisz-psi}
\resultid{ANT-THM-12-03}
\provedhere
لكل \(A>0\)، توجد ثوابت \(c_A,C_A>0\)، غير فعالة في المسار غير المشروط
العام، بحيث لكل \(x\ge3\)، ولكل
\[
1\le q\le(\log x)^A,
\qquad (a,q)=1,
\]
لدينا
\[
\psi(x;q,a)
=
\frac{x}{\varphi(q)}
+O\!\left(C_Axe^{-c_A\sqrt{\log x}}\right),
\tag{12.6}
\]
بانتظام في \(q\) و\(a\).
\end{theorem}

\begin{proof}
نبدأ من مرشح الشخصيات \((12.1)\).

\paragraph{الشخصية الرئيسية.}
إذا كانت \(\chi_0\pmod q\) رئيسية، فهي مستحثة من الشخصية التافهة. من
اللمّة \ref{lem:chapter12-primitive-psi-reduction} ومبرهنة
\ref{thm:chapter12-effective-pnt}:
\[
\psi(x,\chi_0)
=
\psi(x)+O(\log(2q)\log(2x))
=
x+O_A\!\left(xe^{-c\sqrt{\log x}}\right).
\tag{12.7}
\]

\paragraph{الشخصيات غير الرئيسية.}
لتكن \(\chi\pmod q\) غير رئيسية، ولتكن \(\chi^*\pmod r\) جدها البدائي.
تعطي اللمّة \ref{lem:chapter12-nonexceptional-zero-bound} تقديرًا أسيًا
فعالًا لكل الجزء غير الاستثنائي. وإذا وجد صفر استثنائي، تمتص مساهمته
باللمّة \ref{lem:chapter12-exceptional-zero-absorption}. ثم تعيد اللمّة
\ref{lem:chapter12-primitive-psi-reduction} الشخصية المستحثة، فنحصل على
\[
\psi(x,\chi)
\ll_A
xe^{-c_A\sqrt{\log x}}
+
\log(2q)\log(2x).
\]
ولأن \(q\le L^A\)، فإن
\[
\log(2q)\log(2x)
\ll_A L\log(2L)
=o\!\left(xe^{-c_A\sqrt L}\right)
\]
بعد تصغير \(c_A\) عند الحاجة. إذن
\[
\psi(x,\chi)
\ll_A xe^{-c_A\sqrt{\log x}}
\tag{12.8}
\]
موحدًا لكل شخصية غير رئيسية.

ندخل \((12.7)\) و\((12.8)\) في \((12.1)\). تساهم الشخصية الرئيسية بالحد
\(x/\varphi(q)\). أما بقية الشخصيات، فعددها لا يتجاوز \(\varphi(q)\)،
ويلغي عامل \(1/\varphi(q)\) هذه الكثرة من دون خسارة إضافية. تنتج
\((12.6)\).
\end{proof}

\begin{remark}[مقارنة بمسارات أحدث]
توجد صيغ أقوى وأكثر انتظامًا تعزل الصفر الاستثنائي وتستعمل مناطق
Vinogradov--Korobov وتقديرات كثافة خالية من اللوغاريتم وتنافر
Deuring--Heilbronn؛ انظر \cite{ThornerZaman2024}. كما توجد براهين حديثة
من منظور الدوال الضربية الادعائية \cite{Koukoulopoulos2013}. لا يحتاج
المسار الكلاسيكي الأدنى هنا إلى هذه التحسينات.
\end{remark}

\section{الادخار اللوغاريتمي الاعتباطي}

\begin{corollary}[صيغة سيغل--فالفش اللوغاريتمية]
\label{cor:chapter12-siegel-walfisz-log-saving}
\resultid{ANT-COR-12-01}
\provedhere
لكل \(A,B>0\)، يوجد ثابت غير فعال \(C_{A,B}>0\) بحيث، بانتظام في
\(q\le(\log x)^A\) و\((a,q)=1\)،
\[
\psi(x;q,a)
=
\frac{x}{\varphi(q)}
+O_{A,B}\!\left(\frac{x}{(\log x)^B}\right).
\tag{12.9}
\]
\end{corollary}

\begin{proof}
لكل \(c,B>0\)، لدينا
\[
e^{-c\sqrt{\log x}}\ll_{c,B}(\log x)^{-B},
\]
لأن \(\sqrt{\log x}\) يغلب \(\log\log x\). نطبق ذلك على \((12.6)\).
\end{proof}

\section{الانتقال إلى \(\vartheta(x;q,a)\)}

\begin{corollary}[التوزيع المنتظم للدالة \(\vartheta\)]
\label{cor:chapter12-siegel-walfisz-theta}
\resultid{ANT-COR-12-02}
\provedhere
لكل \(A,B>0\)، بانتظام في \(q\le(\log x)^A\) و\((a,q)=1\)،
\[
\vartheta(x;q,a)
=
\frac{x}{\varphi(q)}
+O_{A,B}\!\left(\frac{x}{(\log x)^B}\right),
\tag{12.10}
\]
والثابت غير فعال في الصيغة العامة.
\end{corollary}

\begin{proof}
لدينا
\[
0\le
\psi(x;q,a)-\vartheta(x;q,a)
\le
\sum_{\substack{p^k\le x\\k\ge2}}\log p.
\]
ومن حد القوى الأولية العليا في \textenglish{\texttt{ANT-LEM-09-02}}:
\[
\sum_{\substack{p^k\le x\\k\ge2}}\log p
\ll \sqrt x\log^2(2x).
\]
وهذا أصغر من \(x/(\log x)^B\) لكل \(B\) ثابت عند كبر \(x\). نطبق
\((12.9)\) ونمتص المجال المحدود في الثابت.
\end{proof}

\section{الانتقال إلى \(\pi(x;q,a)\)}

نعرف
\[
\operatorname{Li}(x)=\int_2^x\frac{dt}{\log t}.
\]

\begin{corollary}[التوزيع المنتظم لعدد الأوليات]
\label{cor:chapter12-siegel-walfisz-pi}
\resultid{ANT-COR-12-03}
\provedhere
لكل \(A,B>0\)، بانتظام في \(q\le(\log x)^A\) و\((a,q)=1\)،
\[
\pi(x;q,a)
=
\frac{\operatorname{Li}(x)}{\varphi(q)}
+O_{A,B}\!\left(\frac{x}{(\log x)^B}\right),
\tag{12.11}
\]
والثابت غير فعال في الصيغة العامة.
\end{corollary}

\begin{proof}
من الجمع الجزئي:
\[
\pi(x;q,a)
=
\frac{\vartheta(x;q,a)}{\log x}
+
\int_2^x\frac{\vartheta(t;q,a)}{t\log^2t}\,dt.
\tag{12.12}
\]
لا يجوز تطبيق \((12.10)\) آليًا على كل \(t\le x\)، لأن الشرط
\(q\le(\log t)^A\) قد يفشل في الجزء الصغير. نعالج ذلك بهامش في الأس.

طبق \((12.10)\) بالمعلم \(2A\) وبقوة لوغاريتمية أكبر من \(B+3\)، وضع
\[
y=\exp\!\left(q^{1/(2A)}\right).
\]
لـ\(t\ge y\)، لدينا \(q\le(\log t)^{2A}\)، فتصلح الصيغة الموحدة. ومن
الفرض \(q\le(\log x)^A\) ينتج
\[
y\le e^{\sqrt{\log x}}.
\]

في المجال \([2,y]\)، نستعمل حد تشيبيشيف التافه
\(\vartheta(t;q,a)\le\vartheta(t)\ll t\)، فتكون مساهمته في \((12.12)\)
من رتبة \(O(y)\)، وهي
\[
e^{\sqrt{\log x}}=o\!\left(\frac{x}{(\log x)^B}\right).
\]
وفي المجال \([y,x]\)، نكتب
\[
\vartheta(t;q,a)=\frac{t}{\varphi(q)}+E(t),
\qquad
E(t)\ll_{A,B}\frac{t}{(\log t)^{B+3}}.
\]
تجمع مساهمة الحد الرئيسي في الطرف عند \(x\) والتكامل على \([y,x]\) إلى
\[
\frac1{\varphi(q)}
\left(
\frac{x}{\log x}
+\int_y^x\frac{dt}{\log^2t}
\right)
=
\frac1{\varphi(q)}
\left(
\operatorname{Li}(x)-\operatorname{Li}(y)+\frac{y}{\log y}
\right).
\]
ولذلك تختلف عن \(\operatorname{Li}(x)/\varphi(q)\) بمقدار
\[
O\!\left(\frac{\operatorname{Li}(y)+y/\log y}{\varphi(q)}\right)
=O(y),
\]
وهو ممتص في الخطأ السابق. أما خطأ الطرف والتكامل، فيحققان
\[
\frac{E(x)}{\log x}
+
\int_y^x\frac{E(t)}{t\log^2t}\,dt
\ll_{A,B}
\frac{x}{(\log x)^B}.
\]
وللتفصيل، مساهمة \([y,\sqrt x]\) لا تتجاوز \(O(\sqrt x)\)، ومساهمة
\([\sqrt x,x]\) لا تتجاوز \(O(x/(\log x)^{B+5})\). فتنتج \((12.11)\).
\end{proof}

\section{قراءة النتيجة وحدودها}

\subsection{ما الذي أضيف إلى الفصل العاشر؟}

الفصل العاشر يثبت النهاية عند كل ترديد ثابت، لكنه لا يحدد سرعة التقارب
بانتظام في \(q\). أما هذا الفصل فيعطي مجالًا واحدًا من الترديدات المتغيرة
وحدًا واحدًا صالحًا لها جميعًا. ولذلك يمكن استعمال النتيجة في مسائل تحتاج
جمعًا أو مقارنة على ترديدات لوغاريتمية، وهو ما لا تسمح به النتيجة النوعية
وحدها.

\subsection{لماذا المجال لوغاريتمي؟}

من المنطقة القياسية نحصل، تقريبًا، على عامل
\[
\exp\!\left(-\frac{c\log x}{\log(qT)}\right).
\]
وعند اختيار \(T=e^{\kappa\sqrt{\log x}}\)، يبقى المقام من رتبة
\(\sqrt{\log x}\) ما دام \(\log q\ll\log\log x\)، أي عندما يكون \(q\)
قوة ثابتة من \(\log x\). أما الوصول إلى ترديدات من رتبة قوة موجبة من
\(x\)، فيحتاج آليات أخرى أو متوسطًا على الترديدات.

\subsection{الفعالية وعدم الفعالية}

كل جزء من البرهان قبل اللمّة
\ref{lem:chapter12-exceptional-zero-absorption} يمكن جعله فعالًا ضمن
المداخل المقتبسة الفعالة. عدم الفعالية تدخل فقط من الثابت
\(c_\varepsilon\) في مبرهنة سيغل. وإذا عزل المرء الترديدات المرتبطة
بالشخصية الاستثنائية بدل امتصاصها، أمكن كتابة صيغ فعالة تحتوي حد
\(x^\beta/\beta\) ظاهرًا.

\subsection{ما الذي لا تثبته المبرهنة؟}

لا تثبت سيغل--فالفش:
\begin{itemize}
  \item انتظامًا فرديًا لكل \(q\le x^\delta\) مع \(\delta>0\).
  \item مبرهنة Bombieri--Vinogradov أو مستوى توزيع \(1/2\).
  \item حدًا فعالًا عامًا بعد حذف الصفر الاستثنائي دون تتبعه.
  \item أصغر أولي في فئة حسابية بحد لينك.
  \item أي نتيجة مشروطة بفرضية ريمان المعممة.
\end{itemize}

\section{أمثلة تفسيرية}

\begin{example}[ترديد ثابت]
إذا ثبت \(q=10\)، فإن الفصل العاشر يكفي لإثبات التكافؤ التقاربي. تعطي
سيغل--فالفش أكثر من ذلك: حدًا كميًا صالحًا للترديد نفسه ولجميع الترديدات
حتى قوة ثابتة من \(\log x\).
\end{example}

\begin{example}[ترديد يتحرك ببطء]
إذا كان
\[
q(x)=\left\lfloor(\log x)^3\right\rfloor,
\]
واختيرت فئة \(a(x)\) مختزلة مع \(q(x)\)، فلا يجوز تطبيق مبرهنة الفصل
العاشر مباشرة؛ لأن الترديد لم يعد ثابتًا. أما المبرهنة
\ref{thm:chapter12-siegel-walfisz-psi} فتطبق بالمعلم \(A=3\).
\end{example}

\begin{example}[أثر إشارة الشخصية الاستثنائية]
قبل استعمال سيغل، يظهر الحد الثانوي في الفئة على الصورة
\[
-\frac{\overline{\chi_1(a)}}{\varphi(q)}
\frac{x^{\beta_1}}{\beta_1}.
\]
لذلك يمكن أن تختلف إشارة الانحياز بين الفئات بحسب قيمة الشخصية الحقيقية.
هذه البنية لا تختفي منطقيًا إلا بعد إثبات أن الحد صغير في المجال
اللوغاريتمي. وتبقى آثار الأصفار الاستثنائية مهمة عند دراسة مجالات أكبر أو
متوسطات أدق؛ انظر \cite{DrappeauFiorilli2021,ThornerZaman2024}.
\end{example}

\section{تمارين}

\begin{enumerate}
  \item أثبت بدقة الهوية \((12.1)\) من تعامد شخصيات ديريشليه.
  \item حسن برهان اللمّة \ref{lem:chapter12-primitive-psi-reduction} إلى
  الحد
  \[
  \sum_{p\mid q}\log p\left\lfloor\frac{\log x}{\log p}\right\rfloor.
  \]
  \item تحقق من أن \(L^D e^{-c\sqrt L}\le e^{-c\sqrt L/2}\) لكل \(D\)
  ثابت ولكبر \(L\).
  \item بيّن أن اختيار \(T=e^{\kappa\sqrt L}\) يوازن بين خطأ القطع
  \(x/T\) والمنطقة الخالية القياسية.
  \item حدد موضع استعمال مبرهنة سيغل في برهان المبرهنة المركزية، وفسر لماذا
  لا تنتقل عدم الفعالية إلى اللمّة \ref{lem:chapter12-nonexceptional-zero-bound}.
  \item أثبت أن
  \[
  \sqrt x\log^2x=o\!\left(x(\log x)^{-B}\right)
  \]
  لكل \(B>0\).
  \item أكمل تفاصيل تقدير التكامل في برهان
  \ref{cor:chapter12-siegel-walfisz-pi} بعد تقسيمه عند \(\sqrt x\).
  \item اشرح لماذا لا تكفي مبرهنة الترديد الثابت لإثبات النتيجة عندما
  \(q=\lfloor\log x\rfloor\).
  \item قارن منطقيًا بين عبارة سيغل--فالفش الفردية وعبارة
  Bombieri--Vinogradov المتوسطية، من دون محاولة إثبات الأخيرة.
\end{enumerate}

\section{خلاصة الفصل}

المسار الكامل هو
\[
\text{مرشح الشخصيات}
\longrightarrow
\text{الرد إلى الجد البدائي}
\longrightarrow
\text{الصيغة الصريحة المقطوعة}
\longrightarrow
\text{المنطقة الخالية}
\longrightarrow
\text{عزل الصفر الاستثنائي}
\longrightarrow
\text{مبرهنة سيغل}.
\]
الجزء غير الاستثنائي يعطي ادخارًا أسيًا فعالًا بعد اختيار ارتفاع القطع.
أما الحد الاستثنائي فيمتص فقط باستعمال مبرهنة سيغل، ولذلك يكون الثابت
النهائي غير فعال. وينتج عن ذلك توزيع موحد للأوليات في جميع الفئات المختزلة
للترديدات \(q\le(\log x)^A\)، مع بقاء Bombieri--Vinogradov مرحلة مستقلة
لاحقة.
